\section{Preliminaries}
\subsection{Foundations of Affective Computing}
The goal of affective computing, as articulated by Picard \cite{Picard1997AffectiveComputing}, is to make machines able to identify, understand, and respond to human emotions in a contextually appropriate manner. The primary theoretical challenge in this undertaking is the absence of a universal consensus on the definition and classification of emotion. Two dominant frameworks compete in the literature. Discrete Emotion Theory, prominently advocated by Ekman \cite{Ekman1992Argument}, asserts that a finite set of basic emotions---typically happiness, sadness, anger, fear, disgust, and surprise---are biologically innate and universally recognised across cultures. This view treats emotion recognition as a multi-class classification problem, where each label corresponds to a discrete state with characteristic physiological and behavioural correlates. The Dimensional Emotion Model, associated with Russell \cite{Russell1980Circumplex}, proposes instead that emotions occupy a continuous two-dimensional space defined by valence (pleasant to unpleasant) and arousal (activated to deactivated). While the dimensional model better captures the gradations and blends of real emotional experience, it imposes a substantially more complex regression problem on computational systems.

For this thesis, the EAV corpus's discrete five-class structure (Neutral, Anger, Happiness, Sadness, Calmness) maps naturally to the Ekman-style classification paradigm. We adopt this framework while acknowledging its limitations, particularly its inability to represent the intermediate or ambivalent states that arise during emotion suppression or masking.

A conceptual distinction critical to this thesis is the difference between expressed emotion and experienced emotion. Expressed emotion---reflected in facial muscle movements, vocal prosody, and gestural behaviour---is controlled primarily by the somatic nervous system and is susceptible to voluntary regulation, social masking, and cultural display rules. Experienced emotion---reflected in neurophysiological signals including EEG---is driven by the autonomic nervous system and deep-seated limbic structures that are significantly more difficult to control voluntarily. EEG captures the immediate neurophysiological response to an affective stimulus at the level of synchronised cortical oscillations, providing a measure of the internal state that is less susceptible to the social regulation that governs behavioural expression. The ``affective gap''---the potential discordance between what is expressed and what is experienced---is this thesis's central subject of study.

\subsection{Deep Learning for Feature Extraction}
The effectiveness of any pattern recognition system is constrained by the quality of its feature representation. Early affective computing research relied on handcrafted features: Local Binary Patterns and Histogram of Oriented Gradients for facial expression analysis, and pitch, energy, and Mel-Frequency Cepstral Coefficients for speech emotion recognition. While these methods established important baselines, they frequently lacked robustness to environmental noise and failed to capture the high-level semantic abstractions that deep networks learn automatically.

Deep learning has transformed affective computing from human feature engineering to automated representation learning. For the visual modality, Convolutional Neural Networks (CNNs) and, more recently, Vision Transformers (ViT) have become the dominant architectures. Contemporary work uses pre-trained models fine-tuned on task-specific datasets: EfficientFace and MobileNet for lightweight real-time deployment, and ViT architectures pre-trained on large facial-emotion corpora for higher-accuracy systems at the cost of increased compute. Zhao et al. \cite{Zhao2021LightweightFER} demonstrated that depth-wise separable convolutions as implemented in EfficientFace achieve competitive facial expression recognition accuracy while substantially reducing computational complexity.

For the acoustic modality, the dominant trend has shifted from handcrafted feature engineering to learned representations. He et al. \cite{He2016ResNet} showed that the ResNet architecture, applied to Log-Mel Spectrograms, learns hierarchical acoustic features via residual skip connections while mitigating vanishing gradient issues. The Audio Spectrogram Transformer (AST) \cite{Gong2021AST} further advanced this direction by eliminating the convolutional inductive bias entirely, pre-training a pure attention architecture on the large-scale AudioSet corpus and achieving state-of-the-art results on standard audio classification benchmarks. The AST's pre-training on 2 million diverse audio clips makes it a particularly effective encoder for emotion-relevant audio features, even when fine-tuned on small per-subject datasets.

For EEG, the dominant deep architecture for affective computing tasks is EEGNet \cite{Lawhern2018EEGNet}, a compact depthwise separable CNN designed to process raw EEG signals from arbitrary electrode montages. EEGNet applies a depthwise convolution across the temporal dimension followed by a separable convolution, extracting spatiotemporal patterns in a parameter-efficient manner. Its compact design makes it suitable for per-subject training in the low-data EAV regime, where each subject provides only 280 training trials.

Frequency-domain features are an important complement to raw-signal approaches for EEG. Power Spectral Density (PSD), computed band-by-band across the delta (0.5--4 Hz), theta (4--8 Hz), alpha (8--13 Hz), beta (13--30 Hz) and gamma (30--45 Hz) frequency bands, captures well-established neurophysiological correlates of arousal and valence. Differential Entropy (DE) features compute the logarithm of band-limited power and can be expressed, for a Gaussian band-limited signal with variance $\sigma^2$, as:
\begin{equation}
    h(X) = \frac{1}{2}\log(2\pi e\sigma^2).
\end{equation}
DE features have consistently outperformed PSD-based approaches on standard EEG-emotion benchmarks. On the SEED dataset \cite{Zheng2015SEED}, DE-based methods achieve 7--19 percentage points higher accuracy than PSD when combined with deep classifiers.

\subsection{Multimodal Fusion Strategies}
The performance of a multimodal emotion recognition system depends critically on its fusion strategy---the mechanism by which heterogeneous modality-specific representations are combined into a single prediction. Three broad paradigms are established in the literature.

\textbf{Early fusion} concatenates raw feature vectors from different modalities before classification. This approach is computationally efficient but suffers from the curse of dimensionality and presupposes temporal alignment of audio, visual, and physiological signals---an assumption that is not always met in naturalistic recording conditions. Kalateh et al. \cite{Kalateh2024SystematicReview} report that early fusion frequently performs worse than late fusion on multimodal emotion datasets due to its inability to handle asynchronous temporal dynamics.

\textbf{Late fusion} trains independent classifiers for each modality and aggregates their final probability outputs via voting, averaging, or weighted summation. Late fusion prevents any single modality from dominating the learning process and is robust to missing modalities. Its principal limitation is the loss of cross-modal correlations: it cannot model the relationship whereby a smile changes the interpretation of a sarcastic vocal tone. For this thesis, naive softmax-averaging late fusion serves as the primary comparison baseline, chosen specifically because it represents the simplest possible multimodal combination and any more sophisticated fusion method that fails to beat it carries a negative architectural verdict.

\textbf{Attention-based and hybrid fusion} represents the most recent direction. Inspired by the Transformer architecture \cite{Vaswani2017Attention}, attention mechanisms assign dynamically varying importance weights to each modality, enabling the system to focus on the most reliable signal on a per-sample basis. Cross-modal attention---in which the query from one modality attends to the key-value pairs from another---was formalised for multimodal emotion recognition by Krishna and Patil \cite{KrishnaPatil2020Interspeech} and extended to larger-scale multimodal settings by Lee, Kim and Kim \cite{Lee2024MERCL}, whose architecture employs six directional pairwise attentions across three modalities. For this thesis, a self-attention architecture over a three-token modality sequence is used, which performs equivalently to pairwise cross-attention when the encoders produce single pooled vectors rather than token sequences.

\section{Review of Existing Research}
Soleymani et al. \cite{Soleymani2016EEGFacial} conducted a seminal investigation of continuous emotion recognition by examining synchronisation and cross-modality interference between EEG data and facial expressions. Using the valence-arousal-dominance dimensional model, they employed LSTM-RNN to extract temporal features from EEG power spectral density and facial landmarks. A key contribution was the investigation of whether EEG affective information reflects genuine internal state or is merely an artifact of EMG contamination from facial muscles. Granger causality analysis indicated that physiological signals provide unique information even in contexts without facial cues, establishing the scientific basis for using EEG as an independent channel in multimodal systems.

Xing et al. \cite{Xing2019EEGAudioVisual} examined the interplay between external multimedia stimuli and internal physiological responses by fusing EEG with audio-visual features. Using Differential Entropy features across multiple frequency bands for EEG and a 3D CNN for visual features, they demonstrated that feature-level fusion of EEG with audio-visual features substantially outperformed unimodal baselines on DEAP and MAHNOB-HCI. Their finding that EEG adds a ``hidden layer of emotional context'' that compensates for the ambiguity inherent in external audio-visual cues directly motivates the trimodal approach of the present thesis.

Babu et al. \cite{Babu2024AffectiveFusionNet} proposed AffectiveFusionNet, integrating Vision Transformers with Variational Autoencoders and graph-based relational learning for multimodal emotion recognition. The work demonstrated that global contextual modelling via transformers, combined with structured latent representations, yields superior performance on complex HCI scenarios. This work informs the architectural choice of self-attention over a modality-token sequence in this thesis's fusion module.

Krishna and Patil \cite{KrishnaPatil2020Interspeech} proposed a multimodal emotion recognition framework employing 1D CNNs and cross-modal attention on the IEMOCAP dataset. A key innovation was the formulation of cross-modal attention in which audio features attend to text sequences and vice versa, enabling the model to capture complex inter-modal dependencies not accessible through simple feature concatenation. Their cross-modal attention strategy directly informs the design of the trimodal attention module in this thesis.

Lee, Shomanov, Kabidenova and Yazici \cite{Lee2024EAV} published the EAV dataset---the corpus on which this thesis is evaluated---alongside per-modality baseline results. The dataset comprises 42 subjects in synchronised cue-based recording sessions with 30-channel EEG, audio, and video across five emotion categories. The published per-modality baselines are SCNN audio at 36.7\%, DeepFace vision at 52.8\%, and EEGNet EEG at 36.7\%. The EAV repository provides working data loaders and per-modality trainers but explicitly does not include any fusion implementation, a gap this thesis addresses. A critical structural observation about the EAV corpus is that EEG labels were collected during listening clips and vision labels during speaking clips, reflecting the different behavioural conditions of the cue-based conversation paradigm. This Listen/Speak asymmetry is a known limitation of the corpus that affects the interpretation of cross-modal agreement patterns.

Lee, Kim and Kim \cite{Lee2024MERCL} published the most directly relevant prior work on the methodology this thesis implements. Their framework comprises three stages: modality encoders based on ViT per frame, VGGish audio, and a modified Conformer for EEG, each followed by a residual temporal convolutional network; pre-training with Multimodal Emotion Recognition Contrastive Learning (MERCL), a combination of intra-modal supervised contrastive loss (AMCL), inter-modal contrastive alignment (EMCL), and sample-wise alignment (SMCL); and fine-tuning with six-directional pairwise cross-modal attention fusion. This framework achieves 83.2\% on DEAP and 90.9\% on SEED, but has no published evaluation on EAV and no released code. This thesis implements a simplified variant of the fusion stage on EAV using the EAV repository's pre-existing per-modality encoders, making the methodological contribution a first evaluation of cross-modal attention fusion on the EAV benchmark.

Chumachenko, Iosifidis and Gabbouj \cite{Chumachenko2022SelfAttention} proposed a self-attention fusion framework for audio-visual emotion recognition with incomplete data. Their modality-dropout training scheme, termed ``softhard'' dropout, expands each training batch into sub-batches in which one or more modalities are zeroed, forcing the fusion model to develop predictions robust to single-modality absence. Their framework is directly adapted for the robustness component of this thesis.

Lawhern et al. \cite{Lawhern2018EEGNet} introduced EEGNet, the compact depthwise separable CNN architecture used as the EEG encoder throughout this thesis. EEGNet's design---a depthwise convolution across the temporal axis followed by a pointwise depthwise convolution across channels---enables it to learn spatiotemporal EEG patterns from as few as a few hundred trials per subject, making it the dominant architecture in the low-data EEG affective computing literature.

Gong, Chung and Glass \cite{Gong2021AST} introduced the Audio Spectrogram Transformer (AST), the audio encoder used in this thesis. AST applies a pure Vision Transformer to the mel spectrogram of an audio signal, with patches treated as the token sequence. Pre-trained on AudioSet, which contains 2 million clips and 527 audio event categories, AST achieves state-of-the-art accuracy on AudioSet and transfers effectively to downstream audio classification tasks with relatively small fine-tuning datasets.

Dosovitskiy et al. \cite{Dosovitskiy2021ViT} introduced the Vision Transformer (ViT), demonstrating that a pure attention architecture trained on large image datasets matches or exceeds CNN performance on standard vision benchmarks. The facial-emotion ViT used in this thesis is fine-tuned from a 7-class facial expression model, re-initialised for the EAV 5-class task using the \texttt{ignore\_mismatched\_sizes} loading protocol.

Pan et al. \cite{Pan2024OJEMB} proposed a trimodal framework integrating facial expressions, speech, and EEG across eight emotion categories. Their work, evaluated on a private dataset, reported that EEG integration consistently improves performance over bimodal audio-visual systems, confirming that the EEG modality carries complementary information not captured by external behavioural signals. This provides additional empirical support for including EEG as a third modality despite its lower per-subject accuracy.

Chen et al. \cite{Chen2022EAV} introduced a dataset for emotion recognition with audio, video, EEG and EMG, motivating the broader paradigm of incorporating physiological signals to resolve the ambiguity in audio-visual affect. Their systematic benchmarking of fusion strategies---showing that four-modality fusion outperforms bimodal fusion---supports the trimodal approach of this thesis.

\subsection*{Prior trimodal-fusion methods evaluated on EAV}

Since the EAV dataset's release in September 2024, three trimodal-fusion methods have been published that evaluate on it directly. Each of these is the most relevant prior work for the present thesis and the comparison baselines in Chapter~\ref{chap:within} are drawn from them.

Yin, Shin, Li and Lee \cite{Yin2024AMERL} introduced AMERL, the first attention-based trimodal-fusion method on EAV. AMERL uses modality-tailored transformers as per-modality encoders, then applies self-attention fusion over their pooled features---an architectural design closely related to the present thesis's \texttt{TrimodalAttentionFusion}, although with stronger per-modality encoders and without modality-dropout training. AMERL reports 70.86\% mean accuracy on the 42-subject within-subject protocol. The implementation is publicly released, making AMERL the strongest evaluation-protocol-controlled comparator for our results.

Kang, Li, Gong, Zeng, Yan, Bian, Siok and Wang \cite{Kang2025HyperMML} proposed Hyper-MML, which replaces pairwise attention with a hypergraph fusion module that models higher-order relationships among the three modalities. Hyper-MML reports 76.65\% mean accuracy and macro-$F_1$ 76.80\% on EAV at 42-subject scale, and is described by its authors as the state of the art on the dataset prior to the present work. The code has not been released; we limit our SOTA claim against Hyper-MML accordingly.

EEG-MoCE \cite{Anon2026EEGMoCE} introduces a hyperbolic mixture-of-curvature experts framework, in which each modality is assigned to a Lorentz-manifold expert and the fusion weights are learned in a curvature-aware manner. EEG-MoCE reports 75.88\% mean accuracy on EAV.

None of these three prior methods tests missing-modality robustness, demonstrates a zero-EEG inference path, reports a subject-independent evaluation, or states how its reported epoch or checkpoint was selected. The present thesis addresses all four. Quantitative comparison is presented in Chapter~\ref{chap:within}.

\subsection*{Evaluation protocol, participant variability and selection bias}
\label{sec:lit_evaluation}

Three strands of methodological literature bear directly on the final phase of this thesis.

\textbf{Within-subject and subject-independent evaluation.} Affective computing distinguishes subject-dependent evaluation, in which a model is trained and tested on data from the same person, from subject-independent evaluation, in which the test person is absent from training. The two answer different questions: the first measures how well a system recognises emotion in someone it has already studied, the second how well it generalises to a new user. All three trimodal fusion methods published on EAV report the first. Because physiological and behavioural signals vary strongly between people, the gap between the two protocols is a property of the corpus and the representation, and it has to be measured rather than assumed.

\textbf{Adapting to a new domain through normalisation statistics.} Li et al.\ \cite{Li2016AdaBN} showed that a network trained on one domain can be adapted to another, without labels, by recomputing its batch-normalisation statistics on the target domain. The per-participant standardisation of Chapter~\ref{chap:calibration} applies the same idea at the level of the individual: each participant is treated as a domain, and their feature statistics are re-estimated from a short unlabelled recording.

\textbf{Selection bias and circular analysis.} Choosing a model, a hyperparameter or a stopping epoch on the same data that are then used to report its performance biases the reported figure upward. Varma and Simon \cite{Varma2006CVBias} quantified this bias for cross-validated error estimates, Cawley and Talbot \cite{Cawley2010Overfitting} showed that over-fitting the model-selection criterion can be as damaging as over-fitting the model itself, and Kriegeskorte et al.\ \cite{Kriegeskorte2009DoubleDipping} documented the same ``double dipping'' across neuroscience. Chapter~\ref{chap:audit} measures the size of this bias in the Pre-Thesis~II results.

\textbf{Confidence calibration.} Guo et al.\ \cite{Guo2017Calibration} showed that modern neural networks are often over-confident and that a single temperature parameter fitted on held-out data largely corrects this. Confidence-weighted fusion rules implicitly assume calibrated confidences; Chapter~\ref{chap:calibration} tests that assumption on unseen participants.

\section{Summary of Key Findings}
\subsection{Decoding the EEG--Facial Expression Paradox}
The ``Affective Gap'' refers to the potential mismatch between an individual's subjective internal experience and their observable behaviour. As established by Soleymani et al. \cite{Soleymani2016EEGFacial} and confirmed by subsequent work, EEG provides a more objective and involuntary measure of the internal affective state than facial expressions or vocal prosody, because EEG reflects the direct neurophysiological consequences of limbic and prefrontal processing, without the social-regulatory layer that modulates behavioural expression.

Research on the neurophysiology of emotion suppression \cite{Gross1998EmotionRegulation,Mauss2009Measures} establishes that deliberate suppression of emotional expression is cognitively costly and neurophysiologically detectable. Subjects instructed to suppress visible signs of disgust while watching aversive film clips show increased right prefrontal activation and sustained amygdala activity, both of which produce identifiable EEG signatures, even as their facial expressions remain effectively neutral. This establishes the scientific basis for expecting that a well-trained EEG classifier will sometimes predict a different emotion than the audio-video consensus---not because the EEG model is wrong, but because the EEG signal is tracking a different construct (internal state) than the audio-video system (external expression).

Two primary frequency-domain features are used to extract meaningful information from EEG in affective contexts: Power Spectral Density (PSD) and Differential Entropy (DE). While PSD is the traditional method, DE has emerged as the more potent feature for emotion classification. On the SEED dataset, DE features achieve 7--19 percentage points higher accuracy than PSD when combined with deep learning classifiers. In this thesis, the EEGNet architecture processes raw EEG signals, learning its own internal representation of spatiotemporal patterns rather than relying on a pre-specified frequency-band decomposition.

The principal challenge in EEG-based affective computing is not the signal processing per se but the low-data regime imposed by within-subject EAV protocols. The EAV corpus provides 280 training trials per subject, each of 5 seconds at 100 Hz after downsampling, across a 30-channel electrode montage. This represents a genuinely constrained learning problem: the EEGNet architecture must learn a classifier that generalises to 120 unseen test trials from 280 training examples. The EAV paper's published EEGNet baseline of 36.7\% reflects this difficulty; the present thesis achieves a 43.9\% mean after debugging two substantial training bugs, a 7.2-percentage-point improvement over the published baseline that supports the viability of EEG as a meaningful third modality.

\subsection{The Affective Gap and Its Clinical Significance}
The affective gap is not merely a technical problem for emotion recognition systems; it has direct clinical significance. Several psychological conditions are characterised by systematic discordance between expressed and experienced affect.

\textbf{Alexithymia} is a trait characterised by difficulty identifying, describing, and distinguishing one's own emotional states. Individuals with alexithymia typically exhibit normal or near-normal physiological responses to affective stimuli but have impaired ability to label those responses in the vocabulary of discrete emotions. A multimodal system that separately tracks the EEG signal and the audio-video signal may detect alexithymia as a systematic absence of cross-modal agreement, even when the overall accuracy of both modality-specific classifiers remains intact.

\textbf{Depression and anhedonia} are associated with blunted reward processing in the striatum and diminished beta-band activity in left prefrontal cortex. Depressed individuals may display contextually appropriate facial expressions and vocal prosody, particularly in formal or clinical contexts, while their EEG indicates attenuated positive affect. The discordance between an expressed Happiness and an internally neutral or negative EEG state is precisely the pattern the suppression matrix in this thesis is designed to detect.

\textbf{Autism Spectrum Conditions} are associated with atypical facial expression production and recognition, alongside preserved and sometimes heightened physiological reactivity to social and emotional stimuli. A trimodal system may detect the characteristic pattern of socially-atypical facial expression alongside internally elevated arousal in EEG.

The present thesis establishes the methodology for cross-modal affective coherence analysis but makes no clinical claims, as the EAV corpus consists of healthy participants performing cue-based emotional expressions. The translation of these findings to clinical populations is explicitly deferred to Phase 2 future work.

\subsection{Limitations of Controlled Laboratory Corpora}
The EAV corpus is a cue-based, laboratory-elicited dataset. This design choice enables reliable ground-truth labelling---the emotion category displayed by a participant in a given trial is unambiguous because the participant was instructed to display it---but it introduces a fundamental limitation on the interpretation of cross-modal divergence patterns.

In a cue-based paradigm, participants are instructed to elicit rather than genuinely experience target emotions. The degree of genuine emotional engagement varies between participants and between trials. For participants who produce facial and vocal expressions that match the cue but do not fully enter the target affective state, EEG may indicate a state different from the cue---not because of genuine emotion suppression, but because the emotional simulation was imperfect. This is an alternative explanation for the suppression events detected in our matrix, and the two explanations---genuine suppression versus imperfect elicitation---are not distinguishable from the EAV data alone.

This limitation is explicitly acknowledged in the Results chapter and is the primary motivation for framing clinical validation as Phase 2 work. In a genuinely spontaneous corpus, such as naturalistic conversation recordings with concurrent EEG, the imperfect-elicitation confound disappears, because there is no instructed emotion to simulate. Until such a corpus is available and used for validation, the clinical interpretation of our suppression matrix must remain tentative.

\subsection{Trend and Summary Table}
The development of multimodal emotion recognition shows a clear trend: as models have become more interactive and architecturally sophisticated, the field has moved from treating cross-modal disagreement as noise to be suppressed, toward recognising it as a signal to be characterised. The early feature-concatenation and late-fusion work of the 2010s built the empirical case that multimodal data carries complementary information; the attention-based and contrastive-learning work of the 2020s has developed the architectural tools to extract that information efficiently. The EEG-facial expression paradox---once framed as a source of signal contamination---has been re-conceptualised as an opportunity for affective validation, where internal neural states and external behavioural cues provide complementary rather than redundant views of the emotional state.

\begin{table}[htbp]
\centering
\small
\caption{Research evolution phases in multimodal emotion recognition.}
\label{tab:research_evolution_p2}
\begin{tabular}{p{3.1cm} p{4.0cm} p{3.7cm} p{3.0cm}}
\hline
\textbf{Research evolution phase} & \textbf{Core achievement} & \textbf{Remaining limitation} & \textbf{Representative papers} \\
\hline
Foundational (Physiological Objective) & Established EEG and EMG as objective measures to bypass the Affective Gap. & Severe signal noise; inability to handle dynamic social contexts. & Soleymani et al. \cite{Soleymani2016EEGFacial} \\
\hline
Traditional Fusion (Feature-Level) & Proved that multimodal data through concatenation and stacked autoencoders outperforms unimodal baselines. & Vulnerable to negative transfer; cannot weight modalities based on noise. & Xing et al. \cite{Xing2019EEGAudioVisual}; Pan et al. \cite{Pan2024OJEMB} \\
\hline
Advanced Attention/Hybrid & Introduced dynamic cross-modal interactions through QKV attention and cross-attention to filter and align modalities. & High computational complexity; requires substantial labelled data. & Krishna and Patil \cite{KrishnaPatil2020Interspeech}; Lee et al. \cite{Lee2024MERCL} \\
\hline
Coherence as Signal (This work) & Characterised cross-modal affective incoherence as a primary research target, not a nuisance variable. & Validated on laboratory corpus only; clinical generalisation is Phase 2. & This thesis \\
\hline
\end{tabular}
\end{table}
